\documentclass[10pt]{article}
\usepackage[margin=0.78in]{geometry}
\usepackage{booktabs,tabularx,array,longtable}
\usepackage[T1]{fontenc}
\usepackage{lmodern}
\usepackage{url}
\usepackage{enumitem}
\setlist{nosep,leftmargin=*}
\title{Supplementary File S7: November 1 Simulator Evaluation Evidence}
\author{}
\date{November 1, 2026}
\begin{document}
\maketitle

\section*{Scope and evidence selection}
This file documents the simulator evidence cited in the revised evaluation. It keeps the integrated pipeline, scene--protocol probes, five-seed policy comparison, process-window ablation, approval pilot, and production-mode checks separate because they use different test units and denominators. The raw CSVs, run manifests, reports, and the self-contained visual report are in \texttt{Supplementary-File-S7-data-2026-11-02.zip}. No physical-plant result is included.

\section*{Consolidation update (November 2, 2026)}
The data package is organized into three folders. \texttt{01-paper-evidence/} is the frozen November 1 batch used for every number in the manuscript. \texttt{02-c3-pilot-negative/} contains a November 2 live-Agent multi-round pilot that is \emph{not} used for any claim: after pipeline fixes, the best episode completed with zero governed writes and a process variable outside its band, and a second episode ended at its time budget with zero writes. The pilot exposed platform integration defects (task-artifact addressing, lead queue draining, and goal completion without checking the goal criterion) that remain outside the paper's claims. \texttt{03-system-validation-not-used/} contains integration and demonstration reports that validate the platform but are not part of the controlled experiment blocks. No C3 number appears in the manuscript.

A cast-film probe record from run \texttt{20261101121700-scenario-probes} was excluded. Its manifest declared four scenes but recorded only cast-film, and its simulator endpoint did not match the endpoint in the reproduction command. None of its 24 admissible requests, 24 structural probes, or 20 readbacks are included in the manuscript totals or the data package. The reported scene--protocol counts below use only run \texttt{20261101121221-scenario-probes}.

\section*{Integrated pipeline and scene--protocol probes}
The integrated executions \texttt{20261101103641-9j8} and \texttt{20261101104849-xu0} used seed 42 and reported commit \texttt{da627d21}. Both manifests record that the platform and simulator were already running and that no Agent harness was active. Each run returned 77 passing assertions, zero warnings, and zero failures. Assertions are within-run checks, not independent trial observations.

\begin{table}[h]
\centering
\caption{Valid collected scene--protocol probe outcomes.}
\footnotesize
\setlength{\tabcolsep}{3pt}
\begin{tabularx}{\textwidth}{@{}>{\raggedright\arraybackslash}p{2.0cm}r X r r r@{}}
\toprule
Scene & Writable nodes & Writable nodes by protocol & Admissible accepted & Structural rejected & Readback matched \\
\midrule
Injection & 11 & TCP 5; RTU 2; OPC UA 3; MQTT 1 & 44/44 & 44/44 & 44/44 \\
Wastewater & 7 & TCP 1; RTU 3; OPC UA 2; MQTT 1 & 28/28 & 28/28 & 28/28 \\
Annealing & 7 & TCP 3; RTU 2; OPC UA 1; MQTT 1 & 28/28 & 28/28 & 28/28 \\
Biaxial line & 30 & TCP 20; RTU 1; OPC UA 8; HTTP 1 & 120/120 & 119/119 & 116/116 \\
\bottomrule
\end{tabularx}
\end{table}

TCP and RTU in the table denote Modbus TCP and Modbus RTU. For each writable node, four tested values were expected to pass (three interior values and the upper boundary), followed by four structural-range probes expected to reject. One structural case in the biaxial-line set was not collected. Readback is counted only when the driver returned an actuator value. Across this four-scene subset, the totals are 220/220 accepted admissible requests, 219/219 rejected structural probes, and 216/216 matching available readbacks. The full planned scene--protocol matrix remains incomplete.

The integrated film-line preset check is a separate pipeline result: five endpoint types were configured and nine selected negative probes were rejected (9/9), with no code edits. It reused the same runtime and is not an independent cold-start deployment. The pipeline run manifests and summaries are included in the data archive.

\section*{Five-seed target-band comparison}
C0 is a 60-s no-write condition; its terminal process-variable statistic is the mean of the final five samples. C1 is a deterministic scripted policy that calls the shared governed-write service and uses scenario-specific mission logic; its statistic is the mission-terminal reading. Each seed appears in both conditions for injection mass, wastewater dissolved oxygen (DO), and annealing hardness. The common decision rule checks each recorded statistic against its stated band. C1 per-seed wall time was 150.41--172.90~s, compared with C0's fixed 60~s, so this is a descriptive endpoint contrast rather than a time-matched test. Counts are recomputed from raw \texttt{pv}, \texttt{target}, and \texttt{tol} fields. The runner's \texttt{attained} flag has different condition-specific semantics (C0 combines the PV band with configured guards; C1 records its scenario mission result), so it is retained as source data but not compared across conditions.

\begin{longtable}{@{}p{2.2cm}p{2.5cm}r r r >{\centering\arraybackslash}p{1.8cm}@{}}
\caption{Raw endpoints with a common target-band check recomputed from each process variable, target, and tolerance: 0/15 C0 endpoints and 15/15 C1 endpoints were in band. The condition-specific runner flags are not used for this count.}\label{tab:s7-opt}\\
\toprule
Scene & Target band & Seed & C0 endpoint & C1 endpoint & PV band C0/C1 \\
\midrule
\endfirsthead
\toprule
Scene & Target band & Seed & C0 endpoint & C1 endpoint & PV band C0/C1 \\
\midrule
\endhead
Injection & $32.5\pm0.35$ g & 42 & 31.226 & 32.320 & No / Yes \\
Injection & $32.5\pm0.35$ g & 43 & 31.272 & 32.180 & No / Yes \\
Injection & $32.5\pm0.35$ g & 44 & 31.244 & 32.173 & No / Yes \\
Injection & $32.5\pm0.35$ g & 45 & 31.272 & 32.190 & No / Yes \\
Injection & $32.5\pm0.35$ g & 46 & 31.206 & 32.340 & No / Yes \\
Wastewater DO & $3.4\pm0.6$ mg/L & 42 & 0.202 & 3.864 & No / Yes \\
Wastewater DO & $3.4\pm0.6$ mg/L & 43 & 0.198 & 3.800 & No / Yes \\
Wastewater DO & $3.4\pm0.6$ mg/L & 44 & 0.268 & 3.686 & No / Yes \\
Wastewater DO & $3.4\pm0.6$ mg/L & 45 & 0.200 & 3.576 & No / Yes \\
Wastewater DO & $3.4\pm0.6$ mg/L & 46 & 0.226 & 3.646 & No / Yes \\
Annealing & $95\pm6$ HV & 42 & 134.66 & 97.74 & No / Yes \\
Annealing & $95\pm6$ HV & 43 & 134.60 & 97.44 & No / Yes \\
Annealing & $95\pm6$ HV & 44 & 134.66 & 97.80 & No / Yes \\
Annealing & $95\pm6$ HV & 45 & 134.90 & 99.90 & No / Yes \\
Annealing & $95\pm6$ HV & 46 & 134.36 & 98.20 & No / Yes \\
\bottomrule
\end{longtable}

The individual values above are copied from the ten per-seed CSV result files. Across the five C1 seeds, the scripted missions issued 12 writes in 17 rounds for injection, 56 writes in 59 rounds for wastewater, and 41 writes in 50 rounds for annealing; these are workload counts, not an efficiency comparison. Seed values are simulator initializations, not independent production lines. C1 is not an active-LLM policy and C0 is not an industrial-controller baseline.

\section*{Process-window interlock ablation}
The process-window study used seed 42, three repetitions in each arm, a direct DCW API call, a mock driver, and benchmark mode. The configuration hash was \texttt{3a57e722d4764b98}. Each repetition submitted six negative probes: four hard structural probes and two process-window probes; it also submitted three nominal legal requests and three boundary requests. The additional 20 timed writes per repetition were used for latency summaries and were excluded from these denominators. No readback-mismatch stimulus was included.

\begin{table}[h]
\centering
\caption{Three-repeat process-window ablation. The combined-removal arm removes the process-window and readback-mismatch checks; hard structural limits remain active in every arm.}
\footnotesize
\setlength{\tabcolsep}{3pt}
\begin{tabularx}{\textwidth}{@{}p{3.0cm} >{\centering\arraybackslash}p{2.8cm} >{\centering\arraybackslash}p{4.0cm} >{\centering\arraybackslash}p{2.0cm} >{\centering\arraybackslash}p{2.5cm}@{}}
\toprule
Arm & Negative probes rejected & Window violations reaching driver & Legal false blocks & Boundary expectations met \\
\midrule
Full guard & 18/18 & 0/6 & 0/9 & 9/9 \\
No process-window interlock & 12/18 & 6/6 & 0/9 & 9/9 \\
No-readback check & 18/18 & 0/6 & 0/9 & 9/9 \\
No window + no-readback & 12/18 & 6/6 & 0/9 & 9/9 \\
\bottomrule
\end{tabularx}
\end{table}

The four hard structural probe types were blocked in every arm (12/12 per arm). With the process-window guard active, the two targeted window violations in each repetition did not reach the driver; removing that guard allowed all six across three repetitions to reach the mock driver. The no-readback arm does not test readback-mismatch rejection because that stimulus was absent. Benchmark mode relaxes normal timing and hold-window controls, so these results concern the tested interlock path only.

\section*{Approval, production-mode, and authorization checks}
The approval pilot used one injection scene and a scripted decision fixture. The full CSV and timing records are in the data archive.

\begin{table}[h]
\centering
\caption{Four calls in the approval-route pilot. Auto rows are not clean matched pairs with the manual rows; no operator made a decision.}
\begin{tabularx}{\textwidth}{@{}l X l r r r@{}}
\toprule
Mode & Runner stimulus & Decision & Before & After & Delay (ms) \\
\midrule
Manual & legal-approved & approved & 202 & 204.8 & 10\,593 \\
Manual & legal-rejected & rejected & 216 & 216 & 544 \\
Auto & label: manual arm rejected & n/a & 202 & 204.8 & 10\,071 \\
Auto & label: manual arm approved & n/a & 216 & 218.8 & 10\,073 \\
\bottomrule
\end{tabularx}
\end{table}

The manual reject did not change the simulated control value. The auto-mode label/value mismatch means the four calls do not form a clean paired approval-on/off experiment. This pilot demonstrates route handling only; it is not evidence of operator judgment or an isolated approval-mechanism effect.

In the separate production-mode guard run, the baseline value 190 was accepted. One example each was rejected for interval, step, node-range, recipe-window, and hold-window constraints. The exact API rejection codes are preserved in the source report and CSV. These one-case checks confirm expected branches for their configured inputs; they are not estimates of general rejection rates. A separate three-case authorization probe rejected an unbound write and a cross-node write while accepting one bound control write.

\section*{Reproduction and provenance}
The raw source records in the archive preserve runner outputs and exact run identifiers. Representative commands recorded by the runners are:
\begin{quote}
\footnotesize\ttfamily\raggedright
node bench/pipeline.mjs --profile integrated --seed 42 --cl-seeds 3\par\smallskip
node bench/experiment/scenario-protocol-probes.mjs --scenarios injection,wwtp,anneal,biax --seed 42\par\smallskip
node bench/experiment/optimization-control.mjs --condition C1 --scenarios injection,wwtp,anneal --seed 42\par\smallskip
node bench/experiment/optimization-control.mjs --condition C0 --scenarios injection,wwtp,anneal --seed 42 --duration 60\par\smallskip
node bench/e1-lite.mjs --seed 42 --repeats 3\par\smallskip
node bench/experiment/approval-gate.mjs --scenario injection --base http://127.0.0.1:3005\par\smallskip
node bench/experiment/production-guards.mjs --base http://127.0.0.1:3010\par\smallskip
node bench/experiment/authorization-probe.mjs --scenario injection --base http://127.0.0.1:3005
\end{quote}
The optimization commands were run for seeds 42--46 in both conditions. Re-running them creates new result identifiers and can produce different simulator trajectories if the environment or service state changes. The integrated pipeline reported a commit and harness hash, but the active service source tree was not captured as a hermetic snapshot. The archive includes the run manifests, selected reports, raw CSVs, a machine-readable validation summary, checksum manifest, and self-contained HTML dashboard; it does not include the excluded cast-film run.

\end{document}
